\begin{theorem}[Cyclic trace expansion for a closed MPS chain]\label{thm:trace_eval_word_eq_sum_cyclic}
    \lean{MPSTensor.trace_evalWord_eq_sum_cyclic}
    \leanok
    For a nonempty closed chain, the trace of the word
    $A^{\sigma_0}\cdots A^{\sigma_{L-1}}$ is the sum over cyclic virtual-index
    configurations:
    \begin{align}
        \tr(A^{\sigma_0}\cdots A^{\sigma_{L-1}})
        &=
        \sum_{\alpha_0,\ldots,\alpha_{L-1}}
        \prod_{t=0}^{L-1}
        (A^{\sigma_t})_{\alpha_t,\alpha_{t+1}},
        \qquad \alpha_L=\alpha_0.
        \notag
    \end{align}
\end{theorem}

\begin{proof}
    \leanok
    The trace is the diagonal sum:
    \begin{align}
        \tr(A^{\sigma_0}\cdots A^{\sigma_{L-1}})
        &=
        \sum_{\alpha_0}
        (A^{\sigma_0}\cdots A^{\sigma_{L-1}})_{\alpha_0,\alpha_0}.
        \notag
    \end{align}
    For prescribed endpoints, the matrix product entry expands as
    \begin{align}
        (A^{\sigma_0}\cdots A^{\sigma_{L-1}})_{\alpha_0,\alpha_L}
        &=
        \sum_{\alpha_1,\ldots,\alpha_{L-1}}
        \prod_{t=0}^{L-1}
        (A^{\sigma_t})_{\alpha_t,\alpha_{t+1}}.
        \notag
    \end{align}
    Substituting $\alpha_L=\alpha_0$ in the trace gives the cyclic formula.
\end{proof}

\begin{definition}[Physical isometry in the basic-vector expression]
    \label{def:appendixB_physical_isometry}
    \lean{MPSTensor.AppendixBStructuralData.physicalIsometry}
    \lean{MPSTensor.AppendixBStructuralData.physicalIsometryLeftInverse}
    \leanok
    Suppose that $A$ is equipped with Appendix~B structural data
    $(X,\Lambda,U)$: the matrix $X$ is invertible, the diagonal entries of
    $\Lambda$ are strictly positive, $A^i=X\Lambda U^iX^{-1}$, and $U$
    satisfies the pair-index isometry equation below. Let $\mathcal H_a$ and
    $\mathcal H_b$ have bases indexed by
    $\alpha,\beta\in\{0,\ldots,D-1\}$. The coefficients
    $U^i_{\alpha,\beta}$ define maps
    \begin{align}
        U\colon \mathcal H_a\otimes\mathcal H_b
        &\longrightarrow \mathcal H_{\mathrm{phys}},
        \notag\\
        (Uv)_i
        &=
        \sum_{\alpha,\beta}
        U^i_{\alpha,\beta}v_{\alpha,\beta},
        \notag
    \end{align}
    and
    \begin{align}
        (U^*\psi)_{\alpha,\beta}
        &=
        \sum_i\overline{U^i_{\alpha,\beta}} \psi_i.
        \notag
    \end{align}
\end{definition}

\begin{theorem}[The physical map is an isometric embedding]
    \label{thm:appendixB_physical_isometry_left_inverse}
    \lean{MPSTensor.AppendixBStructuralData.physicalIsometryLeftInverse_comp}
    \lean{MPSTensor.AppendixBStructuralData.physicalIsometry_injective}
    \leanok
    \uses{def:appendixB_physical_isometry}
    Let $A$ be equipped with Appendix~B structural data $(X,\Lambda,U)$ as in
    Definition~\ref{def:appendixB_physical_isometry}. Its pair-index
    isometry equation
    \begin{align}
        \sum_i\overline{U^i_{\alpha,\beta}}
        U^i_{\alpha',\beta'}
        &=
        \delta_{\alpha,\alpha'}\delta_{\beta,\beta'}
        \notag
    \end{align}
    implies $U^*U=1$. In particular, $U$ is injective.
\end{theorem}

\begin{proof}
    \leanok
    Expanding the coefficient of $U^*Uv$ at $(\alpha,\beta)$ and interchanging
    the two finite sums gives
    \begin{align}
        (U^*Uv)_{\alpha,\beta}
        &=
        \sum_{\alpha',\beta'}
        \left(\sum_i\overline{U^i_{\alpha,\beta}}
            U^i_{\alpha',\beta'}\right)
        v_{\alpha',\beta'} =
        v_{\alpha,\beta}.
        \notag
    \end{align}
    A map with a left inverse is injective.
\end{proof}

\begin{definition}[Tensor powers of the physical isometry]
    \label{def:appendixB_physical_isometry_tensor_power}
    \lean{MPSTensor.AppendixBStructuralData.physicalIsometryTensorPower}
    \lean{MPSTensor.AppendixBStructuralData.physicalIsometryTensorPowerLeftInverse}
    \leanok
    \uses{def:appendixB_physical_isometry}
    For every $N\geq 0$, the coefficients of $U$ define
    \begin{align}
        U^{\otimes N}\colon
        (\mathcal H_a\otimes\mathcal H_b)^{\otimes N}
        &\longrightarrow
        \mathcal H_{\mathrm{phys}}^{\otimes N},
        \notag
    \end{align}
    with
    \begin{align}
        (U^{\otimes N}v)_\sigma
        &=
        \sum_{p_0,\ldots,p_{N-1}}
        \left(\prod_{t=0}^{N-1}
            U^{\sigma_t}_{(p_t)_a,(p_t)_b}\right)
        v_{p_0,\ldots,p_{N-1}}.
        \notag
    \end{align}
    The conjugate coefficient map is
    \begin{align}
        (U^{*\otimes N}\psi)_p
        &=
        \sum_\sigma
        \left(\prod_{t=0}^{N-1}
            \overline{U^{\sigma_t}_{(p_t)_a,(p_t)_b}}\right)
        \psi_\sigma.
        \notag
    \end{align}
    This is the tensor power occurring in the basic-vector expression
    \cite[Equation~(3.17)]{Cirac2016MPDO_arXiv}.
\end{definition}

\begin{theorem}[The tensor powers are isometric embeddings]
    \label{thm:appendixB_physical_isometry_tensor_power}
    \lean{MPSTensor.AppendixBStructuralData.physicalIsometryTensorPowerLeftInverse_comp}
    \lean{MPSTensor.AppendixBStructuralData.physicalIsometryTensorPower_injective}
    \leanok
    \uses{def:appendixB_physical_isometry_tensor_power}
    For every $N\geq 0$, one has
    $U^{*\otimes N}U^{\otimes N}=1$. Hence $U^{\otimes N}$ is injective.
\end{theorem}

\begin{proof}
    \leanok
    For virtual-pair configurations $p$ and $q$, interchange the two finite
    sums and factor the sum over physical configurations site by site:
    \begin{align}
        \sum_\sigma
        \prod_t
        \overline{U^{\sigma_t}_{(p_t)_a,(p_t)_b}}
        U^{\sigma_t}_{(q_t)_a,(q_t)_b}
        &=
        \prod_t
        \sum_i
        \overline{U^i_{(p_t)_a,(p_t)_b}}
        U^i_{(q_t)_a,(q_t)_b}.
        \notag
    \end{align}
    The pair-index isometry equation turns the last product into
    $\prod_t\delta_{p_t,q_t}=\delta_{p,q}$. Therefore, the coefficient at $p$
    is $v_p$. A map with a left inverse is injective.
\end{proof}

\begin{theorem}[Appendix B basis change and the parent interaction]
    \label{lem:appendixB_basis_change_parent_interaction}
    \lean{MPSTensor.AppendixBStructuralData.parentInteraction_eq_coreTensor}
    \leanok
    \uses{def:parent_interaction, def:ground_space_parent}
    Let $A^i=X\Lambda U^iX^{-1}$ be the Appendix~B structural form, and write
    $A_{\mathrm{core}}^i=\Lambda U^i$. The canonical local parent interaction
    is unchanged by this basis change:
    $q_L(A)=q_L(A_{\mathrm{core}})$.
\end{theorem}

\begin{proof}
    \leanok
    The basis change gives a gauge equivalence between $A_{\mathrm{core}}$ and
    $A$, hence equality of their length-$L$ local ground spaces. The parent
    interaction is the orthogonal complement projector of this local ground
    space.
\end{proof}

\begin{definition}[Appendix B two-site basic support]
    \label{def:appendixB_two_site_basic_space}
    \lean{MPSTensor.AppendixBStructuralData.twoSiteBasicSpace}
    \leanok
    \uses{def:ground_space_parent}
    For Appendix~B structural data, the two-site basic support is
    $G_2(A_{\mathrm{core}})=G_2(\Lambda U)$.
\end{definition}

\begin{definition}[Two-site bond insertion and basic-vector embedding]
    \label{def:appendixB_two_site_bond_embedding}
    \lean{MPSTensor.AppendixBStructuralData.twoSiteBondInsertion}
    \lean{MPSTensor.AppendixBStructuralData.twoSiteBasicEmbedding}
    \lean{MPSTensor.AppendixBStructuralData.weightedTwoSiteBoundary}
    \leanok
    \uses{def:appendixB_physical_isometry_tensor_power}
    Put $\ket{\varphi}=\sum_b\lambda_b\ket{b,b}$. For outer-boundary
    coefficients $v_{a,c}$, insertion of this vector on the middle virtual
    bond is the map $I_\varphi$ defined by
    \begin{align}
        (I_\varphi v)_{(a,b),(b',c)}
        &=
        \delta_{b,b'}\lambda_bv_{a,c}.
        \notag
    \end{align}
    The two-site physical map is $U^{\otimes2}I_\varphi$. A boundary matrix
    $Y$ determines outer-boundary coefficients by
    $(v_Y)_{a,c}=\lambda_aY_{c,a}$. Opening the outer virtual bond in the
    two-site contraction of
    \cite[Equations~(3.17)--(3.18)]{Cirac2016MPDO_arXiv} gives these
    coefficients.
\end{definition}

\begin{definition}[Virtual bond projector and its adjacent placements]
    \label{def:appendixB_virtual_bond_projectors}
    \lean{MPSTensor.AppendixBStructuralData.virtualBondProjection}
    \lean{MPSTensor.AppendixBStructuralData.leftVirtualBondProjection}
    \lean{MPSTensor.AppendixBStructuralData.rightVirtualBondProjection}
    \leanok
    \uses{def:appendixB_two_site_bond_embedding}
    Write
    $s=\langle\varphi,\varphi\rangle=\sum_b\lambda_b^2$.
    The orthogonal projector onto $\C\varphi$ has pair-index coefficients
    \begin{align}
        (P_\varphi v)_{b,b'}
        &=
        \delta_{b,b'}\frac{\lambda_b}{s}
        \sum_k\lambda_kv_{k,k}.
        \notag
    \end{align}
    On three virtual site pairs, let $P_{01}$ place this projector on the bond
    joining the right index of site zero to the left index of site one, and let
    $P_{12}$ place it on the next bond.
\end{definition}

\begin{theorem}[Commutation of the adjacent virtual bond projectors]
    \label{thm:appendixB_virtual_bond_projectors_commute}
    \lean{MPSTensor.AppendixBStructuralData.leftVirtualBondProjection_comp_right}
    \leanok
    \uses{def:appendixB_virtual_bond_projectors}
    The two adjacent placements commute: $P_{01}P_{12}=P_{12}P_{01}$.
    The following support-location schematic marks the two bonds on which the
    projectors act; it does not depict the doubled operator network for the
    equality above.
    % Schematic formula: U_0--\varphi_{01}--U_1--\varphi_{12}--U_2, with the
    % two virtual-bond supports marked around \varphi_{01} and \varphi_{12}.
    % Ink-to-index: U_r has left/right virtual indices and one upper physical
    % index.  The dots \varphi_{01},\varphi_{12} carry the two virtual
    % half-indices of the corresponding bonds; the red regions locate the two
    % supports and introduce no operator layer or indices.
    % Contractions: the four U--\varphi joins contract the two internal bonds.
    % The two outer virtual indices and all three physical indices remain free.
    % Boundary signature: (virtual-west, virtual-east, physical-up,
    % physical-down)=(1,1,3,0).
    % Source verdict: Equation (3.18) and II_RFP_MPS.png supply the basic-vector
    % geometry.  Lines 1305--1307 only call the implication from RFP structure
    % to a nearest-neighbor commuting parent Hamiltonian immediate; this
    % explicit adjacent-projector lemma and support marking are blueprint
    % algebraic elaborations, not a separately stated source result or figure.
    \begin{center}
        \begin{tenkzfree}
            \tnput[boundary, ports={east:virtual}]
                {bondLeft}{(-12mm,0)}{}
            \tnput[box, ports={west:virtual,east:virtual,north:physical}]
                {bondUzero}{(0,0)}{U_0}
            \tnput[dot, ports={west:virtual,east:virtual},
                label pos=south]
                {bondPhiZero}{(16mm,0)}{\varphi_{01}}
            \tnput[box, ports={west:virtual,east:virtual,north:physical}]
                {bondUone}{(32mm,0)}{U_1}
            \tnput[dot, ports={west:virtual,east:virtual},
                label pos=south]
                {bondPhiOne}{(48mm,0)}{\varphi_{12}}
            \tnput[box, ports={west:virtual,east:virtual,north:physical}]
                {bondUtwo}{(64mm,0)}{U_2}
            \tnput[boundary, ports={west:virtual}]
                {bondRight}{(76mm,0)}{}
            \tnput[boundary, ports={south:physical}]
                {bondPhysicalZero}{(0,14mm)}{}
            \tnput[boundary, ports={south:physical}]
                {bondPhysicalOne}{(32mm,14mm)}{}
            \tnput[boundary, ports={south:physical}]
                {bondPhysicalTwo}{(64mm,14mm)}{}
            \tnjoin[name=bondOuterLeft]{bondLeft.east}{bondUzero.west}
            \tnjoin[name=bondZeroLeft]{bondUzero.east}{bondPhiZero.west}
            \tnjoin[name=bondZeroRight]{bondPhiZero.east}{bondUone.west}
            \tnjoin[name=bondOneLeft]{bondUone.east}{bondPhiOne.west}
            \tnjoin[name=bondOneRight]{bondPhiOne.east}{bondUtwo.west}
            \tnjoin[name=bondOuterRight]{bondUtwo.east}{bondRight.west}
            \tnjoin[name=bondPhysicalJoinZero]
                {bondUzero.north}{bondPhysicalZero.south}
            \tnjoin[name=bondPhysicalJoinOne]
                {bondUone.north}{bondPhysicalOne.south}
            \tnjoin[name=bondPhysicalJoinTwo]
                {bondUtwo.north}{bondPhysicalTwo.south}
            \tnregion[slot=selected, label={$\supp(P_{01})$},
                label pos=south,
                name=bondProjectorZero]{bondPhiZero}
            \tnregion[slot=selected, label={$\supp(P_{12})$},
                label pos=south,
                name=bondProjectorOne]{bondPhiOne}
        \end{tenkzfree}
    \end{center}
\end{theorem}

\begin{proof}
    \leanok
    The contracted bond coordinates are disjoint. Thus replacing the first
    bond by the diagonal basis vector $\ket{k,k}$ and the second by
    $\ket{\ell,\ell}$ is independent of the order. Expanding both rank-one
    projectors gives the same double sum over $(k,\ell)$.
\end{proof}

\begin{theorem}[Surjectivity of the weighted boundary parametrization]
    \label{thm:appendixB_weighted_two_site_boundary_surjective}
    \lean{MPSTensor.AppendixBStructuralData.weightedTwoSiteBoundary_surjective}
    \leanok
    \uses{def:appendixB_two_site_bond_embedding}
    Strict positivity of the $\lambda_a$ makes $Y\mapsto v_Y$ surjective.
\end{theorem}

\begin{proof}
    \leanok
    For prescribed coefficients $v_{a,c}$, take
    $Y_{c,a}=\lambda_a^{-1}v_{a,c}$.
\end{proof}

\begin{theorem}[The two-site basic-vector image]
    \label{thm:appendixB_two_site_basic_embedding_range}
    \lean{MPSTensor.AppendixBStructuralData.twoSiteBasicEmbedding_comp_weightedTwoSiteBoundary}
    \lean{MPSTensor.AppendixBStructuralData.twoSiteBasicEmbedding_range}
    \leanok
    \uses{def:appendixB_two_site_bond_embedding,
        def:appendixB_two_site_basic_space}
    Opening the outer virtual bond in the two-site contraction of
    \cite[Equations~(3.17)--(3.18)]{Cirac2016MPDO_arXiv} gives the following
    local-support family. For every boundary matrix $Y$,
    $U^{\otimes2}I_\varphi(v_Y) = \Gamma_2(\Lambda U)(Y)$.
    Consequently, $\ran(U^{\otimes2}I_\varphi)=G_2(\Lambda U)$.
\end{theorem}

\begin{proof}
    \leanok
    \uses{thm:appendixB_weighted_two_site_boundary_surjective}
    At a two-site physical configuration $(i,j)$, the bond insertion gives
    $\sum_{a,b,c}U^i_{a,b}U^j_{b,c}\lambda_b\lambda_aY_{c,a}$.
    Expanding $\tr((\Lambda U^i)(\Lambda U^j)Y)$ gives the same sum, so
    the two maps agree after the weighted boundary parametrization. Since that
    parametrization is surjective, their ranges are equal.
\end{proof}

\begin{definition}[Two-site basic support projector]
    \label{def:appendixB_two_site_basic_support_projector}
    \lean{MPSTensor.AppendixBStructuralData.twoSiteBasicSupportProjection}
    \leanok
    \uses{def:appendixB_two_site_basic_space}
    Let $P_2$ be the orthogonal projector onto $G_2(\Lambda U)$, as in the
    local parent-space construction of
    \cite[Definition~3.9]{Cirac2016MPDO_arXiv}.
\end{definition}

\begin{theorem}[Range and complement of the two-site basic support projector]
    \label{thm:appendixB_two_site_basic_support_projector}
    \lean{MPSTensor.AppendixBStructuralData.twoSiteBasicSupportProjection_spec}
    \lean{MPSTensor.AppendixBStructuralData.twoSiteBasicEmbedding_range}
    \leanok
    \uses{def:appendixB_two_site_basic_support_projector,
        def:appendixB_two_site_bond_embedding, def:parent_interaction}
    The projector $P_2$ satisfies
    \begin{align}
        P_2^2
        &=
        P_2,
        \notag\\
        \ran(P_2)
        &=
        \ran(U^{\otimes2}I_\varphi) =
        G_2(\Lambda U),
        \notag\\
        P_2
        &=
        1-q_2(\Lambda U).
        \notag
    \end{align}
    In particular, $q_2(\Lambda U)$ annihilates the image of
    $U^{\otimes2}I_\varphi$, while $P_2$ fixes that image pointwise.
\end{theorem}

\begin{proof}
    \leanok
    \uses{thm:appendixB_two_site_basic_embedding_range}
    Orthogonal projection onto $G_2(\Lambda U)$ has that subspace as its range.
    The preceding theorem identifies this range with
    $\ran(U^{\otimes2}I_\varphi)$. Orthogonal projection onto the complementary
    subspace is $q_2(\Lambda U)$, hence $P_2=1-q_2(\Lambda U)$. Since
    $q_2(\Lambda U)^2=q_2(\Lambda U)$,
    \begin{align}
        P_2^2
        &=
        (1-q_2(\Lambda U))^2 =
        1-q_2(\Lambda U) =
        P_2.
        \notag
    \end{align}
    If $\psi$ lies in $\ran(U^{\otimes2}I_\varphi)$, then
    $\psi\in G_2(\Lambda U)=\ker q_2(\Lambda U)$, and therefore
    \begin{align}
        q_2(\Lambda U)\psi
        &=
        0,
        \notag\\
        P_2\psi
        &=
        (1-q_2(\Lambda U))\psi =
        \psi.
        \notag
    \end{align}
\end{proof}

\begin{definition}[Two-site support of a residual isometry family]
    \label{def:residual_family_two_site_support}
    \lean{MPSTensor.residualFamilyPhysicalIsometryMatrix}
    \lean{MPSTensor.ResidualFamilyAppendixBData}
    \lean{MPSTensor.ResidualFamilyAppendixBData.TwoSiteBoundary}
    \lean{MPSTensor.ResidualFamilyAppendixBData.twoSiteBondEmbedding}
    \lean{MPSTensor.ResidualFamilyAppendixBData.twoSiteBondSupportProjection}
    \leanok
    \uses{def:is_residual_isometry_family,
        def:appendixB_two_site_bond_embedding,
        def:appendixB_two_site_basic_support_projector}
    Let
    $\mathcal V=\bigsqcup_j\{j\}\times[D_j]\times[D_j]$
    be the disjoint union of the within-sector virtual pairs, and let
    $V\colon\mathcal V\to\C^d$ be the common physical map whose matrix
    elements are
    $V_{i,(j,\alpha,\beta)}=(U_j^i)_{\alpha,\beta}$.
    For each sector insert the vector
    $\ket{\varphi_j}=\sum_b\lambda_{j,b}\ket{b,b}$ on the middle virtual
    bond, apply $V$ on both sites, and sum the resulting physical vectors.
    This defines the equal-sector two-site embedding
    \begin{align}
        I_{\mathrm{res}}\colon
        \prod_j\C^{D_j\times D_j}
        &\longrightarrow
        (\C^d)^{\otimes2}.
        \notag
    \end{align}
    Unequal-sector virtual bonds do not occur. The corresponding physical
    support operator $P_{2,\mathrm{res}}$ is the orthogonal projector onto
    $G_2(\bigoplus_j B_j)$.
\end{definition}

\begin{theorem}[Range of the residual-family bond embedding]
    \label{thm:residual_family_two_site_support}
    \lean{MPSTensor.IsResidualIsometryFamily.residualFamilyPhysicalIsometryMatrix_isometry}
    \lean{MPSTensor.ResidualFamilyAppendixBData.twoSiteBondEmbedding_range}
    \lean{MPSTensor.ResidualFamilyAppendixBData.twoSiteBondEmbedding_range_eq_directSum_groundSpace}
    \lean{MPSTensor.ResidualFamilyAppendixBData.twoSiteBondSupportProjection_range}
    \lean{MPSTensor.ResidualFamilyAppendixBData.twoSiteBondSupportProjection_eq_complement}
    \lean{MPSTensor.IsBNTCanonicalForm.exists_residualFamilyAppendixBData}
    \leanok
    \uses{thm:residual_isometry_family_one_gram,
        def:residual_family_two_site_support,
        thm:appendixB_two_site_basic_embedding_range,
        def:parent_interaction}
    The common physical map is an isometry, $V^*V=1$. Moreover, the
    equal-sector bond embedding and the physical support projector satisfy
    \begin{align}
        \ran(I_{\mathrm{res}})
        &=
        \sum_j G_2(B_j) =
        G_2\!\left(\bigoplus_j B_j\right),
        \notag\\
        \ran(P_{2,\mathrm{res}})
        &=
        \ran(I_{\mathrm{res}}),
        \notag
    \end{align}
    and
    $P_{2,\mathrm{res}} = 1-q_2\!\left(\bigoplus_j B_j\right)$.
    If the $B_j$ form a basis of normal tensors and their direct sum is a
    renormalization fixed point, the sectorwise Appendix~B decompositions and
    the common residual isometry required above exist.
% Scope restriction (multiplicity-one canonical form): repeated copies remain
% outside this statement; see docs/paper-gaps/cpsv16_rfp_isometry_scope.tex.
\end{theorem}

\begin{proof}
    \leanok
    The first identity is the one-letter Gram identity on the disjoint union
    of the within-sector virtual pairs. In each sector, the bond insertion
    has range $G_2(B_j)$; allowing independent outer-boundary coefficients
    and summing over sectors therefore gives $\sum_jG_2(B_j)$. The local
    ground space of a direct sum is this sum, which proves the two range
    identities. Orthogonal projection onto that space is complementary to
    its parent interaction. Finally, the residual-isometry form of the
    cross-block fixed-point structure supplies the simultaneous sector
    decompositions.
\end{proof}
